\section{Related Work}\label{sec:related}

First, we review the methods for sequential optimisation and parallel search for hyperparameter optimisation that go beyond grid search, random search, and hand tuning. Next, we note that PBT inherits many ideas from genetic algorithms, and so highlight the literature in this space.

The majority of automatic hyperparameter tuning mechanisms are sequential optimisation methods: the result of each training run with a particular set of hyperparameters is used as knowledge to inform the subsequent hyperparameters searched. A lot of previous work uses a Bayesian optimisation framework to incorporate this information by updating the posterior of a Bayesian model of successful hyperparameters for training -- examples include GP-UCB~\citep{srinivas2009gaussian}, TPE~\citep{bergstra2011algorithms}, Spearmint~\citep{snoek2012practical}, and SMAC~\citep{hutter2011sequential}. As noted in the previous section, their sequential nature makes these methods prohibitively slow for expensive optimisation processes. For this reason, a number of methods look to speed up the updating of the posterior with information from each optimisation process by performing early stopping, using intermediate losses to predict final performance, and even modelling the entire time and data dependent optimisation process in order to reduce the number of optimisation steps required for each optimisation process and to better explore the space of promising hyperparameters~\citep{gyorgy2011efficient,agarwal2011oracle,sabharwal2016selecting,swersky2013multi,swersky2014freeze,domhan2015speeding,klein2016fast,snoek2015scalable,springenberg2016bayesian}. 

Many of these Bayesian optimisation approaches can be further sped up by attempting to parallelise them -- training independent models in parallel to quicker update the Bayesian model (though potentially introducing bias)~\citep{shah2015parallel,gonzalez2016batch,wu2016parallel,rasley2017hyperdrive,golovin2017google} -- but these will still require multiple sequential model optimisations. The same problem is also encountered where genetic algorithms are used in place of Bayesian optimisation for hyperparameter evolution~\citep{young2015optimizing}.

Moving away from Bayesian optimisation, even approaches such as Hyperband~\citep{li2016hyperband} which model the problem of hyperparameter selection as a many-armed bandit, and natively incorporate parallelisation, cannot practically be executed within a single training optimisation process due to the large amount of computational resources they would initially require.

The method we present in this paper only requires a single training optimisation process -- it builds upon the unreasonable success of random hyperparameter search which has shown to be very effective~\citep{bergstra2012random}. Random search is effective at finding good regions for sensitive hyperparameters -- PBT identifies these and ensures that these areas are explored more using partially trained models (\eg~by copying their weights). By bootstrapping the evaluation of new hyperparameters during training on partially trained models we eliminate the need for sequential optimisation processes that plagues the methods previously discussed. PBT also allows for adaptive hyperparameters during training, which has been shown, especially for learning rates, to improve optimisation~\citep{loshchilov2016sgdr,smith2017cyclical,masse2015speed}. 

PBT is perhaps most analogous to evolutionary strategies that employ self-adaptive hyperparameter tuning to modify how the genetic algorithm itself operates~\citep{back1998overview,Spears95adaptingcrossover,gloger2004self,clune2008natural} -- in these works, the evolutionary hyperparameters were mutated at a slower rate than the parameters themselves, similar to how hyperparameters in PBT are adjusted at a slower rate than the parameters. In our case however, rather than evolving the parameters of the model, we train them partially with standard optimisation techniques (\eg~gradient descent). Other similar works to ours are Lamarckian evolutionary algorithms in which parameters are inherited whilst hyperparameters are evolved \citep{castillo2006lamarckian} or where the hyperparameters, initial weights, and architecture of networks are evolved, but the parameters are learned \citep{castillo1999g,husken2000optimization}, or where evolution and learning are both applied to the parameters \citep{ku1997exploring,houck1997empirical,igel2005evolutionary}. This mixture of learning and evolutionary-like algorithms has also been explored successfully in other domains~\citep{zhang2011evolutionary} such as neural architecture search~\citep{real2017large,liu2017hierarchical}, feature selection~\citep{xue2016survey}, and parameter learning~\citep{fernando2016convolution}. Finally, there are parallels to work such as \cite{salustowicz1997probabilistic} which perform program search with genetic algorithms.
